% ============================================================
\chapter{DỮ LIỆU VÀ PHƯƠNG PHÁP THỰC NGHIỆM}
\label{ch:du_lieu}
% ============================================================

\section{Bộ dữ liệu}

Báo cáo sử dụng bộ dữ liệu \textbf{Forest Inspection} (mô phỏng bằng AirSim), công bố
trên Zenodo~\cite{forestdataset}. Dữ liệu gồm \textbf{9 chuỗi ảnh} (sequence) thu ở
các tổ hợp góc chụp (pitch $\in\{0^\circ, -60^\circ, -90^\circ\}$) và độ cao bay
($\{30, 50, 80\}$ m), tổng cộng hơn $13.000$ cặp ảnh (RGB + nhãn) ở độ phân giải gốc
$1280\times 768$. Mỗi điểm ảnh được gán một trong \textbf{11 lớp}: Sky,
Deciduous\_trees, Coniferous\_trees, Fallen\_trees, Dirt\_ground, Ground\_vegetation,
Rocks, Building, Fence, Car và Empty (nền/không xác định).

\begin{table}[H]
\centering
\caption{Đặc tả 9 chuỗi ảnh theo góc chụp và độ cao.}
\label{tab:sequences}
\small
\begin{tabular}{c c c c c l}
\toprule
\textbf{Seq} & \textbf{Pitch} & \textbf{Độ cao (m)} & \textbf{Số ảnh} & \textbf{Đặc điểm} \\
\midrule
seq1 & $0^\circ$   & 30 & 1.468 & Ngang, thấp \\
seq2 & $-60^\circ$ & 30 & 1.501 & Chéo, thấp \\
seq3 & $-90^\circ$ & 30 & 1.469 & Thẳng xuống, thấp \\
seq4 & $0^\circ$   & 50 & 1.432 & Ngang, trung bình \\
seq5 & $-60^\circ$ & 50 & 1.447 & Chéo, trung bình \\
seq6 & $-90^\circ$ & 50 & 1.437 & Thẳng xuống, trung bình \\
seq7 & $0^\circ$   & 80 & 1.436 & Ngang, cao \\
seq8 & $-60^\circ$ & 80 & 1.464 & Chéo, cao \\
seq9 & $-90^\circ$ & 80 & 1.397 & Thẳng xuống, cao \\
\bottomrule
\end{tabular}
\end{table}

Quá trình kiểm định bảng màu (\textit{palette}) cho thấy $100\%$ điểm ảnh ánh xạ
chính xác về 11 lớp trên cả 9 chuỗi --- không có điểm ảnh ngoài bảng màu, khẳng định
chất lượng nhãn.

\section{Phân tích mất cân bằng lớp}
\label{sec:imbalance_data}

Bảng~\ref{tab:imbalance} thống kê tần suất điểm ảnh theo lớp (trên tập huấn luyện) và
trọng số nghịch tần suất tương ứng. Chênh lệch giữa lớp lớn nhất
(Ground\_vegetation) và lớp nhỏ nhất (Fence) lên tới \textbf{2.761 lần} --- mức mất
cân bằng rất nghiêm trọng.

\begin{table}[H]
\centering
\caption{Phân bố tần suất lớp và mức mất cân bằng (tập huấn luyện).}
\label{tab:imbalance}
\small
\begin{tabular}{c l r r r}
\toprule
\textbf{Hạng} & \textbf{Lớp} & \textbf{Tần suất (\%)} & \textbf{Tỷ lệ so lớp min} & \textbf{Trọng số (nghịch tần suất)} \\
\midrule
1  & Ground\_vegetation & 32{,}40 & 2.761$\times$ & 0{,}003 \\
2  & Coniferous\_trees  & 27{,}85 & 2.373$\times$ & 0{,}003 \\
3  & Sky                & 20{,}33 & 1.733$\times$ & 0{,}004 \\
4  & Deciduous\_trees   & 15{,}11 & 1.288$\times$ & 0{,}006 \\
5  & Dirt\_ground       & 1{,}83  & 156$\times$   & 0{,}046 \\
6  & Rocks              & 1{,}74  & 148$\times$   & 0{,}049 \\
7  & Fallen\_trees      & 0{,}52  & 44$\times$    & 0{,}164 \\
8  & Building           & 0{,}18  & 15$\times$    & 0{,}478 \\
9  & Car                & 0{,}04  & 3{,}6$\times$ & 2{,}026 \\
10 & \textbf{Fence}     & \textbf{0{,}01} & \textbf{1$\times$} & \textbf{7{,}221} \\
\bottomrule
\end{tabular}
\end{table}

\noindent\textbf{Nhận định:} bốn lớp \textit{Fallen\_trees, Building, Car, Fence}
đều dưới $1\%$; riêng Fence ($0{,}01\%$) và Car ($0{,}04\%$) là \textit{cực hiếm}.
Đây chính là các lớp được kỳ vọng hưởng lợi nhiều nhất từ chiến lược xử lý mất cân
bằng.

\section{Sự thay đổi miền (domain gap)}

Phân bố lớp biến đổi mạnh theo góc chụp. Đáng chú ý, lớp \textbf{Sky chỉ xuất hiện ở
các chuỗi $pitch = 0^\circ$} (chiếm $36$--$58\%$) và \textbf{biến mất hoàn toàn}
($0\%$) ở các chuỗi $-60^\circ$ và $-90^\circ$. Vì vậy, IoU của Sky phụ thuộc hoàn
toàn vào việc tập kiểm thử có chứa chuỗi chụp ngang hay không --- một yếu tố cần lưu
ý khi diễn giải kết quả (Chương~\ref{ch:thao_luan}).

\HINH{Biểu đồ phân bố lớp theo chuỗi/nhóm góc chụp — xuất từ notebook khảo sát dữ liệu.}

\section{Chia tập dữ liệu}
\label{sec:split}

Toàn bộ các thí nghiệm trong báo cáo dùng cách chia theo chuỗi như sau:

\begin{table}[H]
\centering
\caption{Cách chia tập dữ liệu dùng trong các thí nghiệm của báo cáo.}
\label{tab:split}
\begin{tabular}{l l r}
\toprule
\textbf{Tập} & \textbf{Chuỗi} & \textbf{Số ảnh} \\
\midrule
Huấn luyện (Train) & seq1, seq3, seq4, seq5, seq8, seq9 & 8.677 \\
Kiểm định (Val)    & seq2                              & 1.501 \\
Kiểm thử (Test)    & seq6, seq7                        & 2.873 \\
\midrule
\multicolumn{2}{l}{\textbf{Tổng}} & \textbf{13.051} \\
\bottomrule
\end{tabular}
\end{table}

Đây là cách chia \textbf{theo chuỗi} (cross-sequence) chính thức của dự án: chia theo
chuỗi thay vì xáo trộn ngẫu nhiên nhằm \textbf{tránh rò rỉ dữ liệu} giữa các khung
hình liên tiếp của cùng một lần bay, đồng thời buộc tập kiểm thử rơi vào các tổ hợp
góc--độ cao chưa xuất hiện y hệt trong huấn luyện (đánh giá khả năng tổng quát hóa).
Tập kiểm thử \{seq6, seq7\} gồm seq6 ($-90^\circ$, 50 m) và seq7 ($0^\circ$, 80 m).

\noindent\textbf{Lưu ý khi diễn giải.} Vì seq7 chụp ngang ($pitch = 0^\circ$) nên
tập kiểm thử \emph{có} lớp Sky --- đó là lý do IoU của Sky rất cao ($\approx 0{,}99$)
ở mọi thí nghiệm. Điều quan trọng nhất là cả bốn thí nghiệm dùng \textbf{cùng một
cách chia và cùng tập kiểm thử}, nên các con số \textbf{so sánh trực tiếp được với
nhau}. (Một số tài liệu khảo sát ban đầu của dự án từng dùng tập kiểm thử seq9 ---
thẳng xuống, không có Sky; do đó không nên đối chiếu trực tiếp các con số giữa hai
cách chia.)

\section{Kiến trúc và cấu hình huấn luyện}

Để cô lập ảnh hưởng của hàm mất mát, mọi thí nghiệm cố định kiến trúc và siêu tham
số huấn luyện như Bảng~\ref{tab:config}.

\begin{table}[H]
\centering
\caption{Cấu hình huấn luyện dùng chung.}
\label{tab:config}
\begin{tabular}{l l}
\toprule
\textbf{Thành phần} & \textbf{Giá trị} \\
\midrule
Kiến trúc            & UNet++ (\texttt{UnetPlusPlus}) \\
Encoder              & ResNet50, khởi tạo ImageNet \\
Kích thước ảnh       & $1024 \times 1024$ \\
Batch size           & 4 \\
Optimizer            & AdamW (weight decay $10^{-4}$) \\
Learning rate        & $10^{-4}$ \\
Scheduler            & CosineAnnealingLR ($\eta_{min}=10^{-6}$) \\
Mixed precision      & Có (AMP) \\
Số lớp               & 11 (mIoU tính trên 10 lớp, bỏ \textit{Empty}) \\
Tiêu chí chọn mô hình& mIoU cao nhất trên tập kiểm định \\
\bottomrule
\end{tabular}
\end{table}

\subsection{Tăng cường dữ liệu}
Trên tập huấn luyện áp dụng các phép tăng cường (thư viện \texttt{albumentations}):
lật ngang/dọc, xoay $\pm30^\circ$, thay đổi độ sáng--tương phản, biến đổi màu
(hue/saturation), thêm nhiễu Gauss, và chuẩn hóa theo thống kê ImageNet. Tập kiểm
định và kiểm thử \textbf{chỉ} resize và chuẩn hóa (không tăng cường) để đánh giá khách
quan. Tăng cường giúp mô hình bền hơn với biến thiên góc chụp/ánh sáng và giảm
quá khớp --- đặc biệt hữu ích cho các lớp hiếm.

\subsection{Quy trình và mã nguồn}
Toàn bộ quy trình được tổ chức theo bốn notebook: \textbf{NB01} chuẩn bị gói thư viện
và trọng số pretrained offline; \textbf{NB02} huấn luyện và lưu top-$k$ checkpoint
theo mIoU kiểm định; \textbf{NB03} đánh giá trên tập kiểm thử (mIoU, PA, mF1,
per-class IoU, ma trận nhầm lẫn); \textbf{NB04} tổng hợp bảng và hình cho báo cáo. Mã
nguồn dùng chung được module hóa trong thư mục \texttt{src/} (xem
Phụ lục~\ref{ap:reproduce}). Bốn thí nghiệm của báo cáo nằm trong
\texttt{notebooks/executed/}. Luồng dữ liệu giữa bốn notebook minh họa ở
Hình~\ref{fig:pipeline}.

\begin{figure}[H]
\centering
\resizebox{\textwidth}{!}{%
\begin{tikzpicture}[
  >=stealth, font=\footnotesize, node distance=0.5cm,
  nb/.style={draw, rounded corners=3pt, minimum width=2.9cm, minimum height=1.7cm,
             align=center, fill=blue!10},
  art/.style={draw, dashed, rounded corners=2pt, align=center, fill=gray!8,
              font=\scriptsize, minimum height=0.8cm},
  data/.style={draw, cylinder, shape aspect=0.3, aspect=0.3, rotate=90,
               minimum height=2.2cm, minimum width=1.3cm, fill=green!12, align=center,
               font=\scriptsize}]
  % 4 notebook
  \node[nb] (n1) at (0,0)   {\textbf{NB01}\\Chuẩn bị\\assets offline\\{\scriptsize(Net ON, GPU OFF)}};
  \node[nb] (n2) at (4,0)   {\textbf{NB02}\\Huấn luyện\\{\scriptsize(Net OFF, GPU ON)}};
  \node[nb] (n3) at (8,0)   {\textbf{NB03}\\Đánh giá\\{\scriptsize(Net ON, GPU ON)}};
  \node[nb] (n4) at (12,0)  {\textbf{NB04}\\Báo cáo\\{\scriptsize(CPU)}};
  % artefacts giữa các bước
  \node[art, below=0.7cm of n1] (a1) {offline\_assets/\\(gói + weights)};
  \node[art, below=0.7cm of n2] (a2) {checkpoints/\\*.pth};
  \node[art, below=0.7cm of n3] (a3) {results\_all.csv};
  \node[art, below=0.7cm of n4] (a4) {bảng + hình};
  % dataset
  \node[data] (ds) at (6,2.4) {Dataset rừng\\seq1--9};
  % luồng chính
  \draw[->, thick] (n1)--(n2); \draw[->, thick] (n2)--(n3); \draw[->, thick] (n3)--(n4);
  % artefact xuống/lên
  \draw[->] (n1)--(a1); \draw[->] (n2)--(a2); \draw[->] (n3)--(a3); \draw[->] (n4)--(a4);
  \draw[->, gray] (a1) -| (n2); \draw[->, gray] (a2) -| (n3); \draw[->, gray] (a3) -| (n4);
  % dataset -> NB02, NB03
  \draw[->, green!45!black] (ds) -| (n2);
  \draw[->, green!45!black] (ds) -| (n3);
\end{tikzpicture}}
\caption{Pipeline bốn notebook: NB01 chuẩn bị gói và trọng số pretrained offline; NB02
huấn luyện (lưu checkpoint); NB03 đánh giá trên tập kiểm thử (xuất CSV); NB04 tổng hợp
bảng và hình. Dataset rừng cấp dữ liệu cho NB02 và NB03.}
\label{fig:pipeline}
\end{figure}

\section{Thiết kế thí nghiệm}
\label{sec:experiment_design}

Đề tài thực hiện một nghiên cứu \textit{ablation} bốn mức, tăng dần độ phức tạp của
chiến lược xử lý mất cân bằng (Bảng~\ref{tab:experiments}). Cơ sở lý thuyết của từng
hàm mất mát đã trình bày ở Mục~\ref{sec:loss}.

\begin{table}[H]
\centering
\caption{Bốn thí nghiệm so sánh chiến lược xử lý mất cân bằng lớp.}
\label{tab:experiments}
\small
\begin{tabular}{c l l c c}
\toprule
\textbf{TN} & \textbf{Chiến lược} & \textbf{Hàm mất mát} & \textbf{Epochs} & \textbf{Seed} \\
\midrule
Exp 1 & CE thuần (chuẩn đối sánh) & PT.~\eqref{eq:ce} & 50 & 42 \\
Exp 2 & CE + trọng số lớp & PT.~\eqref{eq:ce}--\eqref{eq:enet_w} & 50 & 42 \\
Exp 3 & Focal+Dice+CE (trọng số cố định) & PT.~\eqref{eq:combo} & 50 & 42 \\
Exp 4 & Focal+Dice+CE (trọng số thích nghi) & PT.~\eqref{eq:adaptive} & 100 & 2025 \\
\bottomrule
\end{tabular}
\end{table}

\noindent Các notebook tương ứng đặt trong thư mục \texttt{notebooks/executed/}:
\texttt{exp1\_ce\_only}, \texttt{exp2\_ce\_class\_weights},
\texttt{exp3\_focal\_dice\_ce\_fixed}, \texttt{exp4\_adaptive\_focal\_dice\_seed2025}.

\subsection{Chi tiết cấu hình hàm mất mát}
Để bảo đảm tái lập, cấu hình cụ thể của bốn hàm mất mát (cơ sở lý thuyết ở
Mục~\ref{sec:loss}) như sau:
\begin{itemize}[leftmargin=1.4cm]
  \item \textbf{Exp 1 --- CE thuần:} Cross-Entropy không trọng số, bỏ qua lớp
  \textit{Empty} (\texttt{ignore\_index}).
  \item \textbf{Exp 2 --- CE + trọng số lớp:} dùng PT.~\eqref{eq:enet_w}; trọng số thực
  tế tính trên tập huấn luyện cho ở Bảng~\ref{tab:weights}.
  \item \textbf{Exp 3 --- Focal+Dice+CE (cố định):} Focal ($\gamma=2{,}0$) + Dice (bỏ
  lớp \textit{Empty}) + CE, trọng số cố định $(\lambda_F,\lambda_D,\lambda_{CE})=
  (1{,}0;\,1{,}0;\,0{,}5)$.
  \item \textbf{Exp 4 --- Focal+Dice+CE (thích nghi):} cùng ba thành phần như Exp 3
  nhưng tỷ trọng học được theo PT.~\eqref{eq:adaptive}; khởi tạo $\log\sigma_k=0$
  (tương ứng trọng số ban đầu bằng 1).
\end{itemize}

\begin{table}[H]
\centering
\caption{Trọng số lớp (ENet, $\kappa=1{,}02$) tính trên tập huấn luyện, dùng cho
Exp~2--4.}
\label{tab:weights}
\small
\begin{tabular}{l r r}
\toprule
\textbf{Lớp} & \textbf{Tần suất (\%)} & \textbf{Trọng số $w_c$} \\
\midrule
Sky                & 14{,}33 & 0{,}2375 \\
Deciduous\_trees   & 16{,}44 & 0{,}2123 \\
Coniferous\_trees  & 29{,}04 & 0{,}1329 \\
Fallen\_trees      & 0{,}58  & 1{,}4115 \\
Dirt\_ground       & 2{,}00  & 0{,}9150 \\
Ground\_vegetation & 35{,}46 & 0{,}1129 \\
Rocks              & 1{,}91  & 0{,}9353 \\
Building           & 0{,}19  & 1{,}6549 \\
Fence              & 0{,}013 & 1{,}8018 \\
Car                & 0{,}047 & 1{,}7724 \\
\bottomrule
\end{tabular}
\end{table}
{\footnotesize Trọng số được chuẩn hóa để tổng bằng số lớp; lớp càng hiếm trọng số
càng cao (Fence, Car, Building lớn nhất).}

\section{Môi trường và quy trình}

Huấn luyện thực hiện trên GPU NVIDIA (Kaggle / RTX 6000 PRO). Mỗi lần chạy gồm: thu
thập \& ghép cặp ảnh--nhãn theo chuỗi, tính trọng số lớp (nếu dùng), huấn luyện với
lưu \textit{top-k} checkpoint theo mIoU kiểm định, và đánh giá trên tập kiểm thử bằng
checkpoint tốt nhất.

\TBD{Bổ sung phiên bản thư viện (PyTorch, segmentation-models-pytorch,
albumentations\dots) và thông số phần cứng chính xác nếu trường yêu cầu.}
